\subsection{实验设计与因果推断}

\begin{结论}{配对与区组：用相关性换精度}
两处理比较，每组 $n$：独立样本 $\Var(\bar d)=\frac{2\sigma^2}n$；配对（同一单元前后或匹配对，组内相关 $\rho$）
\[
\重点{\Var(\bar d)=\frac{2\sigma^2(1-\rho)}{n}}.
\]
$\rho=0.8$ 时样本量减为 $\frac15$。区组设计把已知变异来源（批次、日期、个体）作为区组，在区组内随机化，区组效应从误差中扣除。
\end{结论}

\begin{结论}{A/B 测试的样本量“16 法则”}
两比例比较（基线 $p$、检测差 $\Delta$），$\alpha=0.05$ 双侧、检验力 $80\%$：
\[
\重点{n\approx\frac{16\,p(1-p)}{\Delta^2}}\ \text{每组}\quad\bigl(16\approx2(1.96+0.84)^2\bigr).
\]
转化率 $5\%$、检测绝对提升 $0.5\%$：约 $30000$ 人每组。相对效应小、基线低时样本量巨大；方差缩减（CUPED：用实验前指标回归调整）可大幅降低需求。
\end{结论}

\begin{结论}{$2^k$ 全因子设计}
$k$ 个两水平因子、$N=2^k$ 次运行：主效应 $=\bar y_+-\bar y_-$，交互作用用对应列乘积的正负号分组；每个效应的方差 \重点{\frac{4\sigma^2}{N}}，所有运行都参与每个效应的估计（比“一次改一个因子”效率高 $k$ 倍量级，且能估计交互）。
无重复时用 Daniel 正态概率图找显著效应；中心点检验弯曲。
\end{结论}

\begin{结论}{部分因子设计与混杂}
$2^{k-p}$ 设计：用生成元定义 $p$ 个因子（如 $D=ABC$），代价是效应混杂（别名）。分辨率 R：
III 型主效应与二阶交互混杂；IV 型主效应不与二阶交互混杂；\重点{\text{V 型二阶交互之间也不混杂}}。$2^{7-4}_{\text{III}}$ 用 $8$ 次运行筛选 $7$ 个因子（饱和设计）；Plackett--Burman 设计 $N=12,20,\dots$。
\end{结论}

\begin{结论}{拉丁方：同时控制两个干扰因素}
$n\times n$ 拉丁方：行、列为两个区组因素，字母为处理，每处理在每行每列恰出现一次。$n^2$ 次运行估计 $n$ 个处理效应，误差自由度 $(n-1)(n-2)$。
正交拉丁方（希腊--拉丁方）可再加一个因素；\重点{n=6\text{ 时不存在}}（Euler 三十六军官问题，Tarry 1900），其余 $n\ge3$ 都存在。
\end{结论}

\begin{结论}{最优设计：把点放在哪里}
线性回归 $y=a+bx$，$x\in[-1,1]$，$N$ 个点：$\Var\hat b=\frac{\sigma^2}{\sum(x_i-\bar x)^2}$，最优是\重点{\text{一半点放 }-1\text{、一半放 }+1}（$\Var\hat b=\frac{\sigma^2}N$）。
D 最优：最大化 $\det(X^{\mathsf T}X)$；二次模型的 D 最优在 $\{-1,0,1\}$ 各放 $\frac13$。代价：无法检验线性假设——实践中加中心点与少量中间点。
\end{结论}

\begin{结论}{响应面方法}
先用一阶设计沿最速上升方向移动（步长 $\propto$ 系数），接近最优时用二阶设计拟合 $y=\beta_0+\mathbf x^{\mathsf T}\mathbf b+\mathbf x^{\mathsf T}B\mathbf x$：驻点 \重点{\mathbf x_s=-\frac12B^{-1}\mathbf b}，$B$ 的特征值符号判断极大、极小或鞍点（典型分析）。
中心复合设计：$2^k$ 角点 $+$ $2k$ 轴点（距离 $\alpha=(2^k)^{1/4}$ 时可旋转）$+$ 中心点。
\end{结论}

\begin{结论}{序贯监测：偷看数据的代价}
固定 $\alpha=0.05$ 在每次中期分析都检验：$5$ 次检验的实际一类错误约 $14\%$，无限次偷看趋于 $1$（重对数律）。
\重点{\text{alpha 消耗函数}}：O'Brien--Fleming 早期界很严（$z\approx4.56$ 于 $\frac15$ 信息量）、最终界接近 $1.96$；Pocock 各次相同界（$5$ 次时 $z\approx2.41$）。始终有效的方法：序贯似然比、置信序列。
\end{结论}

\begin{结论}{效应量与检验力的经验尺度}
Cohen's $d=\frac{\mu_1-\mu_2}{\sigma}$：$0.2$ 小、$0.5$ 中、$0.8$ 大。两样本 $t$ 检验 $80\%$ 检验力每组需要 $n\approx\frac{16}{d^2}$：$d=0.5$ 要 $64$ 人，$d=0.2$ 要 $400$ 人。
低检验力研究的显著结果高估效应（“赢家诅咒”/Type M 误差）：\重点{\text{检验力 }20\%\text{ 时显著结果的效应平均被夸大约 }2\text{ 倍以上}}。
\end{结论}

\begin{结论}{混杂、对撞与后门准则}
因果图中，调整变量集 $Z$ 能识别 $X\to Y$ 的效应 $\iff$ $Z$ 阻断所有从 $X$ 指向“进入 $X$”的后门路径且不含 $X$ 的后代。
\重点{\text{不要调整对撞变量}}（$X\to C\leftarrow Y$）：条件化会制造虚假相关（Berkson 悖论：只看住院病人时两种疾病负相关；只看录取者时成绩与面试负相关）。中介变量调整后得到的是直接效应。
\end{结论}

\begin{结论}{准实验设计：DID 与断点回归}
双重差分：$\hat\tau=(\bar y_{T,\text{后}}-\bar y_{T,\text{前}})-(\bar y_{C,\text{后}}-\bar y_{C,\text{前}})$，需平行趋势假设（检查处理前趋势）。
断点回归：处理由连续变量过阈值决定，比较阈值两侧的回归极限，\重点{\text{局部随机化}}——只识别阈值附近的效应；带宽选择、操纵检验（McCrary 密度检验）是关键。
\end{结论}

\begin{结论}{倾向得分与逆概率加权}
倾向得分 $e(x)=P(T=1\mid x)$：在可忽略性（无未测混杂）下，以 $e(x)$ 为条件即平衡全部协变量。IPW 估计平均处理效应
\[
\重点{\hat\tau=\frac1n\sum\Bigl[\frac{T_iY_i}{e(x_i)}-\frac{(1-T_i)Y_i}{1-e(x_i)}\Bigr]},
\]
$e$ 接近 $0$ 或 $1$ 时权重爆炸（截断或改用匹配/双重稳健估计）。倾向得分不能修正未测混杂。
\end{结论}

\begin{结论}{整群随机化的设计效应}
每群 $m$ 人、组内相关 $\mathrm{ICC}=\rho$：方差膨胀因子
\[
\重点{\mathrm{DE}=1+(m-1)\rho},
\]
有效样本量 $\frac{N}{\mathrm{DE}}$。$m=50$、$\rho=0.05$ 时 $\mathrm{DE}=3.45$。增加群数比增加群内人数有效得多；分析时必须考虑聚类（否则标准误被低估）。
\end{结论}

\begin{结论}{测量信度：Spearman--Brown 公式}
单次测量信度 $\rho$（真分数方差占比），平均 $k$ 次平行测量的信度
\[
\重点{\rho_k=\frac{k\rho}{1+(k-1)\rho}}.
\]
信度 $0.5$ 的测验加长到 $4$ 倍信度 $0.8$。相关系数的衰减修正 $r_{\text{真}}=\frac{r_{xy}}{\sqrt{\rho_x\rho_y}}$——测量噪声使相关被低估（与回归中的衰减偏差同源）。
\end{结论}

\begin{结论}{分层抽样：Neyman 分配}
总体分 $H$ 层（规模 $N_h$、层内标准差 $\sigma_h$），总样本 $n$：方差最小的分配
\[
\重点{n_h\propto N_h\sigma_h},
\]
（考虑单位成本 $c_h$ 时 $n_h\propto\frac{N_h\sigma_h}{\sqrt{c_h}}$）。层间差异越大、层内越同质，分层收益越大；比例分配 $n_h\propto N_h$ 永不比简单随机抽样差。
\end{结论}

\参考{Box--Hunter--Hunter, \emph{Statistics for Experimenters}；Montgomery, \emph{Design and Analysis of Experiments}；Pearl, \emph{Causality}；Hernán--Robins, \emph{Causal Inference: What If}；Kohavi 等, \emph{Trustworthy Online Controlled Experiments}；MIT 14.387、15.077。}
